\chapter{Összefoglalás}

A dolgozat célja egy olyan C++ könyvtár megtervezése és megvalósítása volt,
amely Linux rendszereken segíti a folyamatok operációs rendszer szintű
sandboxolását. A munka középpontjában két kernelmechanizmus állt: a seccomp,
amely a folyamat által használható rendszerhívások szűrésére alkalmas, valamint
a Landlock, amely fájlrendszer-hozzáférési szabályok beállítását teszi
lehetővé.

A koncepció és a tervezés során a dolgozat bemutatta a folyamat-szintű
sandboxolás célját, valamint azokat a Linux-specifikus technológiákat, amelyekre
a könyvtár épült. Fontos szempont volt, hogy a megoldás ne rejtse el teljesen a
kernelmechanizmusok korlátait, de a közvetlen rendszerhívásokhoz,
BPF-utasításokhoz és Landlock bitmaszkokhoz képest mégis kényelmesebb,
magasabb szintű C++ felületet adjon.

A megvalósítás részeként elkészült egy seccomp-alapú szabályépítő modul. Ez
előre definiált kategóriák alapján gyűjti össze az engedélyezendő
rendszerhívások számait, majd ezekből a véglegesítéskor BPF szűrőt állít elő.
Az egymást követő rendszerhívásszámok tartományszabályként is bekerülhetnek a
szűrőbe, így a létrejövő program tömörebb maradhat. A rendszerhívások
meghívásához a könyvtár saját, alacsony szintű wrapper függvényeket használ,
amelyek közvetlenül a Linux rendszerhívási interfészre épülnek.

A Landlock modulban elkészült a \code{LandlockRuleSet} osztály, amely RAII
jelleggel kezeli a Landlock ruleset fájlleírót. A modul magasabb szintű
hozzáférési kategóriákat vezet be, például olvasási, írási és végrehajtási
jogosultságokat, majd ezeket alakítja át a kernel által elvárt bitmaszkokká. A
szabályok hozzáadásakor az implementáció figyelembe veszi, hogy a célpont fájl
vagy könyvtár-e, és ennek megfelelően szűkíti az alkalmazott hozzáférési
maszkot.

A könyvtár működését GoogleTest alapú tesztek és demonstrációs példaprogramok
ellenőrizték. A tesztek nemcsak az inicializálás és a szabályok létrehozásának
sikerességét vizsgálták, hanem azt is, hogy a korlátozások ténylegesen
érvényesülnek-e. A Landlock tesztek engedélyezett és tiltott
fájlrendszer-hozzáféréseket vizsgáltak, míg a seccomp tesztek ellenőrizték,
hogy az engedélyezett rendszerhívások végrehajthatók maradnak, a tiltottak
pedig a folyamat megszakítását eredményezik. A \code{cat} és
\code{http\_\allowbreak server} példaprogramok azt mutatták be, hogy a két
mechanizmus egyszerűbb programhelyzetekben együtt is alkalmazható.

A dolgozat ismert sérülékenységek alapján azt is bemutatta, hogy a seccomp és a
Landlock nem a sérülékenységek kiváltó okát szüntetik meg, hanem a sikeres
kihasználás hatását csökkenthetik. Az ImageMagick és az Apache HTTP Server
példái alapján látható, hogy egy nem megbízható bemenetet feldolgozó folyamat
esetén hasznos lehet a fájlrendszer-hozzáférések és a rendszerhívások előzetes
szűkítése. Ez a többrétegű védekezés elvéhez illeszkedik: a sérülékeny
szoftvert továbbra is javítani kell, de a folyamat jogosultságainak
korlátozása önálló védelmi réteget adhat.

A megvalósítás jelenlegi formájában elsősorban demonstrációs és kutatási célú.
További fejlesztési lehetőség lenne több előre definiált seccomp kategória
kidolgozása, akár részben automatikusan, például a
\code{systemd-analyze syscall-filter} által használt rendszerhívás-csoportok
felhasználásával. Hasznos lenne a könyvtár egyszerűbb csomagolása és
telepíthetővé tétele is, hogy könnyebben integrálható legyen más C++
projektekbe. Emellett egy már létező alkalmazásba történő beépítés pontosabban
megmutathatná, hogy gyakorlati környezetben milyen szabályprofilokra és API-beli
bővítésekre lenne szükség.

Összességében a dolgozat bemutatta, hogy a seccomp és a Landlock alkalmasak
lehetnek Linux folyamatok jogosultságainak csökkentésére, és ezek a
mechanizmusok egy magasabb szintű C++ könyvtáron keresztül kezelhetőbb formában
is elérhetővé tehetők. A megvalósított könyvtár nem teljes körű sandbox-rendszer,
de szemlélteti, hogyan építhető fejlesztőbarátabb API a Linux által biztosított
alacsony szintű biztonsági primitívekre.
